\titleframe

% Contoh Kode Slide
%\section{Introduction}
%\begin{frame}{\insertsectionhead}
%	\framesubtitle{A short introduction to Trigon}
%	\themename is a modern, elegant and versatile theme for Beamer, inspired by
%	the
%	\href{https://github.com/matze/mtheme}{\textsc{metropolis} theme} from Matthias
%	Vogelgesang.
%	\vfill
%	\themename comes with lots of nice extra features
%	\begin{itemize}
%		\item Multiple style variations for title, section and normal slides
%		\item Simple customization of theme colors
%		\item Lots of convenient options to tweak the design
%	\end{itemize}
%\end{frame}

% Contoh Kode Tabel
%  \begin{table}[H]
% \centering
% \caption{A nice table example}
% \begin{tabular}{@{} lccc @{}}
%    \toprule
%    & \textbf{Velocity} & \textbf{Angle}  & \textbf{Vertical force} \\
%    & $U$ & $\alpha$  & $F_z$ \\
%    & [m/s] & [$^\circ$]  & [N] \\
%    \midrule
%    2D simulation  & 9 & 2 & 9.23 \\
%    3D simulation  & 10.0 & 3 & 15.039 \\
%    Experiment A   & 11.31 & 2.5 & 13.2 \\
%    Experiment B   & 11.26 & 2.7 & 12.6 \\
%    Experiment C   & 11.33 & 2.47 & 13.6 \\
%    \bottomrule
%  \end{tabular}
%\end{table}

% Contoh Kode Gambar
%  \begin{figure}[ht!]
%	\begin{subfigure}[b]{0.3\textwidth}
%		\frame{\includegraphics[width=\textwidth]{layout_example-03.jpg}}
%		\caption*{plain}
%	\end{subfigure}
%	\hspace{\fill}
%	\begin{subfigure}[b]{0.3\textwidth}
%		\frame{\includegraphics[width=\textwidth]{layout_example-02.jpg}}
%		\caption*{style1}
%	\end{subfigure}
%	\hspace{\fill}
%	\begin{subfigure}[b]{0.3\textwidth}
%		\frame{\includegraphics[width=\textwidth]{layout_example-01.jpg}}
%		\caption*{style2 (default)}
%	\end{subfigure}
%\end{figure}

% Contoh Kode Blocks
% \begin{block}{Regular block}
%	Just a regular block
%\end{block}
%\begin{alertblock}{Alert block}
%	Some important thing
%\end{alertblock}
%\begin{exampleblock}{Example block}
%	No difference with regular block to avoid excessive distraction
%\end{exampleblock}

\setbeamertemplate{frame footer}{Alauddin Maulana Hirzan, M. Kom - Universitas Semarang}

\section{Keamanan dan Privasi \textit{Internet of Things} \#2}
\begin{frame}{\insertsectionhead}
	\framesubtitle{Apa itu Privasi?}
	\justifying
	Definisi Privasi dibagi menjadi dua:
	\vfill
	\begin{itemize}
		\item Kebebasan untuk menyendiri tanpa ada gangguan bisa disebut dengan privasi
		\item Kebebasan dalam mengendalikan informasi personal atau pribadi 
	\end{itemize}
	\vfill
	Namun dalam ranah IoT, yang menjadi titik berat adalah privasi dalam hal informasi.
	\vfill
	\begin{block}{Info}
		Informasi seperti umur, status pernikahan, tempat tinggal merupakan salah satu bentuk informasi pribadi yang seharusnya tidak disebar.
	\end{block}
\end{frame}

\begin{frame}{\insertsectionhead}
	\framesubtitle{Data Personal \#1}
	\justifying
	Menurut \textbf{Undang-Undang Nomor 24 Tahun 2013 tentang Perubahan Atas Undang-Undang Nomor 23 Tahun 2006 tentang Administrasi Kependudukan}, dalam pasal 1 angka 22 menyatakan bahwa Data Pribadi adalah data perseorangan tertentu yang disimpan, dirawat dan dijaga kebenaran serta dilindungi kerahasiaannya.
	\vfill
	Namun sayangnya di Indonesia tidak menjelaskan secara detail data apa saja yang bisa dianggap sebagai data pribadi.
	\vfill
	\begin{block}{Info}
		Ditambah dengan awamnya sebagian besar rakyat Indonesia akan pentingnya data pribadi, memperparah keadaan yang ada.
	\end{block}
\end{frame}

\begin{frame}{\insertsectionhead}
	\framesubtitle{Data Personal \#2}
	\justifying
	Menurut \textbf{\textit{Office of the Australian Information Commissioner}}, mendefinisikan data pribadi sebagai berikut:
	\vfill
	\begin{itemize}
		\item 'informasi sensitif' (termasuk informasi atau pendapat tentang asal ras atau etnis seseorang, pendapat politik, keyakinan agama, catatan kriminal)
		\item 'informasi kesehatan' (yang juga merupakan 'informasi sensitif')
		\item 'informasi kartu kredit'
		\item informasi 'catatan karyawan' (dengan pengecualian), dan
		\item 'informasi nomor file pajak'
	\end{itemize}
\end{frame}

\begin{frame}{\insertsectionhead}
	\framesubtitle{Data Personal \#3}
	\justifying
	Contoh umum mengenai informasi pribadi:
	\vfill
	\begin{itemize}
		\item Informasi tentang kehidupan pribadi atau keluarga seseorang : Nama seseorang, tanda tangan, alamat rumah, alamat email, nomor telepon, tanggal lahir, catatan medis, detail rekening bank
		Informasi tentang kebiasaan dan praktik kerja seseorang.
		\item Detail pekerjaan seseorang : seperti alamat kerja dan detail kontak, gaji, jabatan, dan praktik kerja.
	\end{itemize}
	\vfill
	\begin{block}{Info}
		Menanyakan pekerjaan yang dilakukan sekaligus gaji itu melanggar privasi, apalagi dilakukan di luar negera.
	\end{block}
\end{frame}

\begin{frame}{\insertsectionhead}
	\framesubtitle{Mengapa Privasi Penting? \#1}
	\justifying
	Dalam masyarakat modern, pembahasan seputar privasi adalah perdebatan tentang kebebasan modern. Berikut ini adalah alasan mengapa privasi itu penting:
	\vfill
	\begin{itemize}
		\item 1. Membatasi Kekuatan Negara : Privasi adalah batas kekuasaan pemerintah, serta kekuasaan perusahaan swasta. Semakin seseorang mengetahui tentang kita, semakin besar kekuatan yang mereka miliki atas kita.
		\item 2. Menghormati Individu : Privasi adalah tentang menghormati individu. Jika seseorang memiliki keinginan yang masuk akal untuk merahasiakan sesuatu, mengabaikan keinginan orang tersebut tanpa alasan yang kuat adalah tindakan yang tidak sopan.
	\end{itemize}
\end{frame}

\begin{frame}{\insertsectionhead}
	\framesubtitle{Mengapa Privasi Penting? \#2}
	\justifying
	\begin{itemize}
		\item 3. Manajemen Reputasi : Privasi memungkinkan orang untuk mengelola reputasi mereka. Bagaimana kita dinilai oleh orang lain memengaruhi peluang, persahabatan, dan kesejahteraan kita secara keseluruhan.
		\item 4. Menjaga Batas Sosial yang Tepat: Orang-orang menetapkan batas-batas dari orang lain dalam masyarakat. Batas-batas ini bersifat fisik dan informasi. Kita membutuhkan tempat menyendiri untuk mundur, tempat di mana kita bebas dari pandangan orang lain untuk bersantai dan merasa nyaman.
		\item 5. Kepercayaan; Dalam hubungan, baik pribadi, profesional, pemerintahan, atau komersial, kita bergantung pada kepercayaan pihak lain. Pelanggaran kerahasiaan adalah pelanggaran kepercayaan itu.
	\end{itemize}
\end{frame}

\begin{frame}{\insertsectionhead}
	\framesubtitle{Mengapa Privasi Penting? \#3}
	\justifying
	\begin{itemize}
		\item 6. Kontrol Atas Kehidupan Seseorang : Data pribadi sangat penting untuk begitu banyak keputusan yang dibuat tentang kita, mulai dari apakah kita mendapatkan pinjaman, lisensi, atau pekerjaan hingga reputasi pribadi dan profesional kita.
	\end{itemize}
	\vfill
	\begin{block}{Info}
		Semakin kita banyak mengekspos diri sendiri ke sumber terbuka (Internet), maka ancaman yang bisa terjadi akan semakin tinggi.
	\end{block}
\end{frame}

\begin{frame}{\insertsectionhead}
	\framesubtitle{Mengapa Privasi Penting? \#4}
	\justifying
	 \begin{figure}[ht!]
			\begin{subfigure}[b]{0.85\textwidth}
					\frame{\includegraphics[width=\textwidth]{china.png}}
				\end{subfigure}
		\end{figure}
\end{frame}

\begin{frame}{\insertsectionhead}
	\framesubtitle{Mengapa Privasi Penting? \#5}
	\justifying
	\begin{figure}[ht!]
		\begin{subfigure}[b]{0.63\textwidth}
			\frame{\includegraphics[width=\textwidth]{china2.png}}
		\end{subfigure}
	\end{figure}
\end{frame}

\begin{frame}{\insertsectionhead}
	\framesubtitle{Privasi di Dunia}
	\justifying
	\begin{figure}[ht!]
		\begin{subfigure}[b]{0.5\textwidth}
			\frame{\includegraphics[width=\textwidth]{privasi-dunia.png}}
		\end{subfigure}
	\end{figure}
	\vfill
	\begin{block}{Info}
		Per tahun 2022, Indonesia berada di peringkat 59 dengan skor privasi hanya mencapai 39.9.
	\end{block}
\end{frame}

\begin{frame}{\insertsectionhead}
	\framesubtitle{Privasi dan \textit{Internet of Things} - Mata-mata}
	\justifying
	Dengan adanya teknologi IoT ini, membuat perangkat-perangkat semakin dekat dan menjangkau sebagian besar daerah yang ada di Indonesia. Namun tentunya hal ini memiliki nilai positif dan negatif bagi kita semua sebagai pengguna teknologi tersebut. Jika dalam nilai positif, perangkat IoT dapat membantu pekerjaan manusia. Tetapi dalam nilai-nilai negatif juga memberikan kerugian, yaitu:
	\vfill
	\begin{itemize}
		\item Terganggunya privasi  pengguna yang dapat dilakukan oleh sebagian perangkat-perangkat IoT jika perangkat tersebut diambil alih fungsi oleh \textbf{pihak ketiga}
	\end{itemize}
\end{frame}

\begin{frame}{\insertsectionhead}
	\framesubtitle{Privasi dan \textit{Internet of Things} - Data Melimpah}
	\justifying
	Data yang dihasilkan oleh perangkat IoT oleh kurang dari 10.000 rumah tangga dapat menghasilkan hampir 150 juta data diskrit setiap hari. Data diskrit merupakan data yang bersifat berkesinambungan sebagai contoh adalah perintah ke AC.
	\vfill
	Oleh karena itu, pengguna dapat dengan jelas melihat peningkatan kemungkinan pelanggaran privasi di IoT. Pengguna memiliki lebih banyak titik masuk bagi peretas sambil membiarkan informasi sensitif dan perangkat IoT pengguna tetap rentan.
\end{frame}

\begin{frame}{\insertsectionhead}
	\framesubtitle{Privasi dan \textit{Internet of Things} - Paparan Publik yang Tidak Diinginkan}
	\justifying
	Berikutnya yang paling penting di antara masalah privasi di IoT adalah paparan publik yang tidak diinginkan. Hal ini tertulis jelas di dokumentasi produsen yang panjang tersebut, namun pengguna sering kali tidak membacanya.
	\vfill
	Misalnya, perusahaan asuransi dapat mengumpulkan informasi dari individu mengenai kebiasaan mengemudi mereka melalui mobil yang terhubung.
	\vfill
	\begin{block}{Info}
		Cara ini yang digunakan oleh sosial media agar mereka tetap bisa memberikan akses gratis ke semua orang. Yaitu \textbf{menjual data pribadi mereka}
	\end{block}
\end{frame}	

\begin{frame}{\insertsectionhead}
	\framesubtitle{\textit{Case Study} - Sisi Gelap \textit{Voice Assistant}}
	\justifying
	\begin{figure}[ht!]
		\begin{subfigure}[b]{0.8\textwidth}
			\frame{\includegraphics[width=\textwidth]{voice.jpg}}
		\end{subfigure}
	\end{figure}
\end{frame}

\begin{frame}{\insertsectionhead}
	\framesubtitle{\textit{Case Study} - \textit{Amazon Alexa} \#1}
	\justifying
	\begin{figure}[ht!]
		\begin{subfigure}[b]{0.7\textwidth}
			\frame{\includegraphics[width=\textwidth]{alexa.jpg}}
		\end{subfigure}
	\end{figure}
\end{frame}

\begin{frame}{\insertsectionhead}
	\framesubtitle{\textit{Case Study} - \textit{Amazon Alexa} \#2}
	\justifying
	Dikutip dari \textbf{avast.com}:
	\vfill
	\begin{quote}
		Kekhawatiran privasi Amazon Alexa berlimpah. Perangkat speaker pintar Echo yang telah terhubung memiliki beberapa kesalahan besar, termasuk secara tidak sengaja merekam percakapan dan mengirimkannya ke kolega. Amazon juga secara tidak sengaja mengirim keseluruhan rekaman suara Echo dari satu pengguna Jerman ke orang yang salah.
	\end{quote}
	\vfill
	Apabila disalahgunakan, informasi-informasi tersebut menjadi senjata yang digunakan ke pemilik utama
\end{frame}

\begin{frame}{\insertsectionhead}
	\framesubtitle{Privasi dan \textit{Internet of Things} - \textit{What to do?}}
	\justifying
	Jika kita menggunakan layanan IoT seperti \textit{Alexa} maupun \textit{Google Nest}. Maka perhatikan pertanyaan-pertanyaan berikut:
	\vfill
	\begin{itemize}
		\item Apakah perangkat IoT saya mengumpulkan data pribadi? Jika ya, apa jenis datanya?
		Ke mana data dikirim?
		\item Bagaimana data ini digunakan? Siapa yang memiliki akses ke data?
		\item Di mana data disimpan? Berapa lama data ini disimpan?
		\item Apakah data aman selama transfer?
		\item Bagaimana pelanggaran data akan ditangani dan entitas mana yang akan diberi tahu jika terjadi pelanggaran data?
	\end{itemize}
\end{frame}

\begin{frame}{\insertsectionhead}
	\framesubtitle{Privasi dan \textit{Internet of Things} - \textit{What to do?}}
	\justifying
	Ada beberapa hal yang bisa dilakukan untuk meningkatkan keamanan privasi dari teknologi IoT, seperti:
	\vfill
	\begin{itemize}
		\item \textit{Self-Host}. Takut data diambil perusahaan besar untuk dijual? Buat sendiri sistem IoT yang dapat dikontrol.
		\item Kontrol data di layanan yang digunakan.
		\item Keamanan dasar adalah senjata utama
		\item  VPN mengamankan komunikasi data antara klien dan server
	\end{itemize}
\end{frame}

\setbeamertemplate{background} 
{
	\includegraphics[width=\paperwidth,height=\paperheight]{protect.jpg}
}
\begin{frame}[plain]
\end{frame}

\setbeamertemplate{background} 
{
	\includegraphics[width=\paperwidth,height=\paperheight]{thank-you.jpg}
}
\begin{frame}[plain]
\end{frame}